\subsection{Comparisons at fixed degree}\label{sec:degree-history}
Classical degree bounds of order \(\Delta^2\log\Delta\) use successive replications and births of new cycles \cite{ChristopherLloyd1995,LiChanChung,HanLi2012}. Han and Li obtain degree bounds at every sufficiently large integer and sharper estimates on specified dyadic sequences. The Li\'enard degree bounds of Xiong and Han, the odd-multiplicity cycles studied by Chang, Zhang and Han, and the canard cycles of \'{A}lvarez, Coll, De Maesschalck and Prohens also concern degree or component degrees \cite{XiongHan2014,ChangZhangHan2013,AlvarezEtAl2020}. Our Table~\ref{tab:fold} requires a count of expanded monomials and simultaneous hyperbolicity under a common cap.

A two-source renewal preprint reports \(\Delta^2\log_2\Delta-O(\Delta^2)\) at unrestricted support \cite{LuRenewal}. Its Taylor interpolation matches a prescribed finite jet at a point where the congruence divisor is nonzero, while Lemma~\ref{lem:whole-disk-division} realizes the integral variation for each admissible numerator by canceling the apparent poles of \(A/D_0\) throughout an oval disk whose vertical projection may contain zeros of \(D_0\).

Qin and Sun's October preprint constructs hyperbolic cycles on an anisotropic Chebyshev grid, using two filtrations of mixed coefficients \cite[Theorem~1.1 and \S\S3.2, 6.2--6.3]{QinSun2026}. A mechanical-renewal preprint combines outer births, reflection and fourfold copying \cite[Proposition~3.3]{LuMechanical}. Both squared-logarithmic bounds have order \(\Delta^2(\log\Delta)^2\) in the degree. At their respective scales \(n\), the sufficient term ceilings \(n^2+1\) and \(2n^2\) for the displayed perturbations give order \(N(\log N)^2\) under the diagonal budget. The minimum support needed for those recurrences is open. The ordinary support calculations for the displayed fields are in \S\S\othernumber{sup:renewal-support}{S1.1}--\othernumber{sup:mechanical-budget}{S1.2} of the electronic supplement \cite{LuElectronic}.

For iterated pullbacks \((u,v)\mapsto(p(u),p(v))\) with \(\deg p\ge2\), a preprint gives a bound of \(O(\Delta^2)\) in the resulting degree \(\Delta\) when each step only lifts existing cycles through full branches from a fixed seed \cite[Theorem~2]{EshkobilovEtAl2026}. In the present construction the regular saddle coefficients and the critical returns let one univariate drive serve whole clusters. Varying the potential supplies further coefficients without introducing mixed powers. The resulting exact-degree bound in Corollary~\ref{cor:degree-only} uses only \(\lceil3\Delta/2\rceil\) terms.
